% -*- mode: latex; coding: utf-8 -*-
%-----------------------------------------------------------------------------
% P3CTeX -- Manual d'usuari de la classe
% Part of P3CTeX. Author: Scvria.
%
% Build from directory tex/ (TEXINPUTS so class and packages are found):
%   set TEXINPUTS=.;./latex;./code;
%   pdflatex -interaction=nonstopmode doc/P3CTeX.tex
%-----------------------------------------------------------------------------
\documentclass[CAT]{P3CTeX}

\begin{document}

\title{La classe \texttt{P3CTeX}\\
  \large Plantilla per a PAC, PEC i projectes UOC}
\author{Scvria\\P3CTeX / UOC PEC Repository}
\date{Versió 0.5 -- 2026/02}
\maketitle

\tableofcontents

\section{Introducció}

La classe \texttt{P3CTeX} proporciona una plantilla cohesionada per a
treballs de l'estudiantat de la UOC (PAC, PEC, pràctiques de projectes,
etc.). La classe:
\begin{itemize}
  \item configura el format de pàgina, tipografia i encapçalaments
        d'acord amb l'estil UOC;
  \item integra les biblioteques del projecte P3CTeX
        (\texttt{pxCORE}, \texttt{pxPRP}, \texttt{pxUML},
        \texttt{pxSRC}, \texttt{pxGDX}, \dots);
  \item ofereix una interfície d'usuari \LaTeX2e senzilla basada en
        opcions de classe i unes poques macros d'alt nivell.
\end{itemize}

Aquest document descriu:
\begin{itemize}
  \item les opcions de classe (família de claus \texttt{P3CTeX});
  \item les claus de metadades acadèmiques (\texttt{P3CTeX/data});
  \item la relació amb el paquet \texttt{pxCORE} i com passar opcions
        a les biblioteques independents;
  \item les macros principals per treballar amb la portada i el sumari.
\end{itemize}

\section{Carregar la classe}

L'ús bàsic és:
\begin{verbatim}
\documentclass[<opcions>]{P3CTeX}
\end{verbatim}

On \verb|<opcions>| és una llista separada per comes d'opcions de la
família de claus \texttt{P3CTeX} (vegeu la Secció~\ref{sec:class-opts}),
possiblement combinades amb opcions de \texttt{pxCORE} (Secció
\ref{sec:pxcore-bridge}).

Per defecte, si no s'indica cap idioma, la classe treballa en català.

\section{Opcions de classe P3CTeX}\label{sec:class-opts}

Les opcions de classe es defineixen internament a la família de claus
\verb|P3CTeX|. Es poden passar directament a \verb|\documentclass| com
en l'exemple següent:
\begin{verbatim}
\documentclass[PR,EN]{P3CTeX}
\end{verbatim}

\subsection*{Control de portada, sumari i declaració de plagi}

\begin{description}
  \item[\texttt{cover} (bool)] controla si s'imprimeix la portada
    institucional. Per defecte la clau existeix amb valor inicial
    \texttt{false}, però el valor per defecte quan no s'especifica és
    \texttt{true} (és a dir, \texttt{cover} està pensat per ser activat per
    les meta-claus com \texttt{PR}).
  \item[\texttt{no-cover}] meta-clau equivalent a \texttt{cover=\{false\}}.
  \item[\texttt{toc} (bool)] controla si es genera el sumari (\emph{table
    of contents}). Per defecte inicialment \texttt{false}, però amb
    valor per defecte \texttt{true}.
  \item[\texttt{no-toc}] meta-clau equivalent a \texttt{toc=\{false\}}.
  \item[\texttt{no-plagi} (bool)] controla si es mostra la declaració de
    no plagi a la portada o a l'inici del sumari. El valor inicial és
    \texttt{false} i el valor per defecte \texttt{true}; les meta-claus
    definides més avall l'activen automàticament.
\end{description}

\subsection*{Meta-profils per a tipus de lliurament}

Per comoditat, la classe ofereix diverses meta-claus que combinen
opcions de portada, sumari, declaració de plagi i idioma:

\begin{description}
  \item[\texttt{CA}] perfil per a \emph{Control d'Avaluació} en anglès:
    desactiva la portada, activa sumari i declaració de no plagi, i
    selecciona idioma anglès (equivalent a
    \verb|cover=false,toc=true,no-plagi=true,en|).
  \item[\texttt{PAC}] perfil per a \emph{Prova d'Avaluació Contínua} en
    català (usa \verb|cat| com a idioma).
  \item[\texttt{PEC}] perfil per a \emph{Prova d'Avaluació
    Continuada/avaluació} en castellà (usa \verb|es| com a idioma).
  \item[\texttt{PR}] perfil per a memòries o projectes amb portada
    completa i sumari activat (equivalent a
    \verb|cover=true,toc=true|).
  \item[\texttt{help}] perfil de ``guia'': activa portada, sumari,
    declaració de no plagi i idioma anglès, i carrega dades d'exemple
    (vegeu la família \verb|P3CTeX/data|).
  \item[\texttt{fake}] perfil de demostració amb dades d'exemple
    (nom d'estudiant, assignatura, etc.) i idioma anglès; útil per
    generar plantilles o exemples sense dades reals.
\end{description}

\subsection*{Selecció d'idioma}

La classe ofereix tres idiomes principals, implementats via
\texttt{babel} i una configuració interna de \texttt{pxCORE}:

\begin{description}
  \item[\texttt{EN} / \texttt{en}] carrega \verb|\usepackage[english]{babel}|
    i activa la configuració interna en anglès.
  \item[\texttt{ES} / \texttt{es}] carrega \verb|\usepackage[spanish]{babel}|.
  \item[\texttt{CAT} / \texttt{cat}] carrega \verb|\usepackage[catalan]{babel}|. Si cap
    opció d'idioma no s'especifica, aquest és l'efecte per defecte.
\end{description}

Les formes en majúscules (\verb|EN|, \verb|ES|, \verb|CAT|) són
meta-claus que simplement apliquen la clau corresponent en minúscula.

\subsection*{Altres opcions}

\begin{description}
  \item[\texttt{data}] permet passar un conjunt de claus de la família
    \texttt{P3CTeX/data} (vegeu la Secció~\ref{sec:data}). Exemple:
\begin{verbatim}
\documentclass[
  PR,
  data = {
    student       = {Nom Cognom},
    subj-fullname = {Nom complet de l'assignatura},
    ca-fullname   = {PAC 1: Títol},
  }
]{P3CTeX}
\end{verbatim}
  \item[\texttt{debug} (bool)] activa la traça de depuració compartida
    amb \texttt{pxCORE} (vegeu \ref{sec:pxcore-bridge}); pensat per a
    desenvolupament de la plantilla.
  \item[\texttt{print-doc} (bool)] reserva un indicador per a futurs
    fluxos de ``documentació impresa''; actualment s'emmagatzema però
    no canvia el comportament.
\end{description}

\section{Metadades acadèmiques}\label{sec:data}

A més de les opcions de comportament, la classe manté una petita
estructura de dades amb la informació acadèmica bàsica. Aquestes dades
s'utilitzen per omplir la portada, encapçalaments i alguns textos
interns.

Les claus de \verb|P3CTeX/data| es passen mitjançant l'opció de classe
\verb|data={...}|:

\begin{description}
  \item[\texttt{student}] nom complet de l'estudiant.
        Valor per defecte: \texttt{Nom\textasciitilde Alumne}.
  \item[\texttt{grade}] nom del grau o màster
        (per exemple, \texttt{Enginyeria\textasciitilde Informàtica}).
  \item[\texttt{subj-fullname}] nom complet de l'assignatura.
  \item[\texttt{subj-shortname}] nom curt de l'assignatura (per a
        encapçalaments).
  \item[\texttt{subj-code}] codi de l'assignatura (per exemple, \texttt{01312}).
  \item[\texttt{ca-fullname}] nom complet de l'activitat (PAC, PEC,
        pràctica, projecte, \dots).
  \item[\texttt{ca-shortname}] nom curt de l'activitat (PAC1, PR, etc.).
  \item[\texttt{date}] data de lliurament; per defecte \verb|\today|.
\end{description}

Addicionalment hi ha dues meta-claus útils per a proves:

\begin{description}
  \item[\texttt{example}] emplena totes les claus anteriors amb dades
    d'exemple realistes (nom de l'estudiant, codi d'assignatura, etc.).
  \item[\texttt{help}] estableix cada clau amb el seu propi nom (per
    exemple, \texttt{student=student}) per visualitzar fàcilment on es
    fa servir cada dada a la portada i encapçalaments.
\end{description}

\section{Passar opcions a pxCORE}\label{sec:pxcore-bridge}

La classe \texttt{P3CTeX} no carrega directament les biblioteques
independents (\texttt{pxUML}, \texttt{pxSRC}, \dots); en lloc d'això,
carrega el paquet d'agregació \texttt{pxCORE} amb una cadena
d'opcions.

Tota opció de classe que \texttt{P3CTeX} no reconeix explícitament es
rea\-grupa en una llista interna i es passa a \texttt{pxCORE} com a
opcions de paquet. Això permet escriure:
\begin{verbatim}
\documentclass[
  PR,
  UML,
  SRC,
  UMLconfig = {maincolor=blue!70!black}
]{P3CTeX}
\end{verbatim}

On:
\begin{itemize}
  \item \verb|UML| i \verb|SRC| són claus de la família \texttt{pxCORE}
        que activen respectivament les biblioteques \texttt{pxUML} i
        \texttt{pxSRC};
  \item \verb|UMLconfig| és una clau de \texttt{pxCORE} que passa
        opcions directament a \texttt{pxUML}.
\end{itemize}

Per a una descripció completa de les claus de \texttt{pxCORE}, vegeu el
manual \verb|pxCORE.tex|.

\section{Macros d'alt nivell de P3CTeX}

Un cop carregada la classe, algunes macros d'alt nivell simplifiquen el
treball amb la portada i el sumari (\texttt{pxGDX}) i amb l'estil UOC.
En aquesta secció es recull la seva sintaxi bàsica, l'ús típic i un
exemple curt.

\subsection{PECTeXconfig}

\paragraph*{Sinopsi}
\begin{verbatim}
\PECTeXconfig{<opcions>}
\end{verbatim}

\paragraph*{Ús típic}
Aplica un conjunt d'opcions de configuració al mòdul gràfic intern
\texttt{pxGDX}. Les opcions es reparteixen habitualment entre tres blocs:
\begin{itemize}
  \item \texttt{cover-config}: claus de \texttt{gdxCOVER} per a la portada;
  \item \texttt{page-config}: claus de \texttt{gdxSTYLE} per a pàgina;
  \item \texttt{pdf-config}: claus de \texttt{gdxPDF} per a metadades.
\end{itemize}
Per al catàleg complet de claus, vegeu el manual \texttt{pxGDX.tex}.

\paragraph*{Exemple}
\begin{verbatim}
\PECTeXconfig{
  cover-config = {
    text-backtitle = {Grau en Enginyeria Informàtica},
    text-title     = {Nom de l'assignatura},
    text-minititle = {PAC1},
    text-author    = {Nom Cognom}
  },
  page-config = {
    author-name = {Nom Cognom},
    doc-name    = {PAC 1},
    subj-name   = {Nom Curt},
    refcode     = {01312}
  }
}
\end{verbatim}

\subsection{PECTeXprintCover}

\paragraph*{Sinopsi}
\begin{verbatim}
\PECTeXprintCover[*]
\end{verbatim}

\paragraph*{Ús típic}
Imprimeix la portada configurada a través de \texttt{pxGDX}. Sense
estrella, utilitza les opcions de portada vigents. La forma amb estrella
(\texttt{*}) força determinades opcions internes (com ara \texttt{noplagi})
per tal que la declaració de no plagi es gestioni des de la portada.

En fluxos normals no cal cridar-la manualment: la classe la invoca al
començament del document segons l'estat de les opcions \texttt{cover} i
\texttt{no-plagi}.

\paragraph*{Exemple}
\begin{verbatim}
% Forçar la impressió explícita de la portada
\PECTeXprintCover*
\end{verbatim}

\subsection{PECTeXprintTOC}

\paragraph*{Sinopsi}
\begin{verbatim}
\PECTeXprintTOC[*]
\end{verbatim}

\paragraph*{Ús típic}
Imprimeix el sumari del document. Sense estrella, mostra un sumari
estàndard. Amb estrella (\texttt{*}), afegeix abans del sumari la
declaració de no plagi quan aquesta no s'ha imprès a la portada.

També en aquest cas, la macro es crida automàticament si l'opció
\texttt{toc} és activa, de manera que la majoria de documents no la
necessiten en el cos principal.

\paragraph*{Exemple}
\begin{verbatim}
% Imprimir només el sumari, amb declaració de no plagi prèvia
\PECTeXprintTOC*
\end{verbatim}

\subsection{Macros d'estil per al cos del document}

Les macros següents, definides per la classe sobre el mòdul
\texttt{pxGDX}, ajuden a mantenir una aparença coherent en el cos del
document.

\subsubsection{sectionuoc i variants}

\paragraph*{Sinopsi}
\begin{verbatim}
\sectionuoc{<títol>}
\subsectionuoc{<títol>}
\subsubsectionuoc{<títol>}
\end{verbatim}

\paragraph*{Ús típic}
Aquestes macros creen seccions, subseccions i subsubseccions amb la
paleta de colors corporatius definida per \texttt{pxGDX}, mantenint
coherent l'estil visual dels enunciats i seccions principals.

\paragraph*{Exemple}
\begin{verbatim}
\sectionuoc{Introducció}
Text de la secció introductòria.
\end{verbatim}

\subsubsection{pxHRule}

\paragraph*{Sinopsi}
\begin{verbatim}
\pxHRule
\pxHRule*
\pxHRule[<amplada>]
\pxHRule*[<amplada>]
\end{verbatim}

\paragraph*{Ús típic}
Traça una línia horitzontal utilitzada com a separador. La forma amb
estrella pot variar lleugerament el gruix o estil intern; l'argument
opcional permet limitar l'amplada de la línia.

\paragraph*{Exemple}
\begin{verbatim}
Enunciat de l'exercici.
\pxHRule[0.5\linewidth]
Text de la resposta...
\end{verbatim}

\subsubsection{pxPageRule}

\paragraph*{Sinopsi}
\begin{verbatim}
\pxPageRule
\pxPageRule*
\end{verbatim}

\paragraph*{Ús típic}
Dibuixa una línia horitzontal centrada a la pàgina, pensada per separar
blocs dins d'enunciats llargs o entre seccions d'un mateix exercici.

\paragraph*{Exemple}
\begin{verbatim}
Apartat A ...

\pxPageRule*

Apartat B ...
\end{verbatim}

\section{Patrons d'ús recomanats}

Tot i que les opcions de \texttt{P3CTeX} es poden combinar de moltes
maneres, en la pràctica hi ha alguns patrons que es repeteixen sovint.
En aquesta secció se'n presenten dos: una PAC típica i una memòria de
projecte (perfil \texttt{PR}).

\subsection{PAC típica}

Per a una PAC estàndard en català és habitual:
\begin{itemize}
  \item utilitzar el perfil \texttt{PAC} com a base (\texttt{cover} desactivat,
        sumari i declaració de no plagi activats, idioma català);
  \item omplir totes les dades acadèmiques via \texttt{data=\{\dots\}};
  \item ajustar la configuració de \texttt{pxGDX} mínimament (títol curt,
        nom de l'estudiant i codi d'assignatura).
\end{itemize}

Un esquelet típic seria:
\begin{verbatim}
\documentclass[
  PAC,
  data = {
    student       = {Nom Cognom},
    grade         = {Enginyeria Informàtica},
    subj-fullname = {Nom complet de l'assignatura},
    subj-shortname= {NCA},
    subj-code     = {01312},
    ca-fullname   = {PAC 1: Títol},
    ca-shortname  = {PAC1}
  }
]{P3CTeX}

\begin{document}

\PECTeXconfig{
  page-config = {
    author-name = {Nom Cognom},
    doc-name    = {PAC 1},
    subj-name   = {NCA},
    refcode     = {01312}
  }
}

% Portada desactivada; sumari i declaració de no plagi segons perfil PAC.

\sectionuoc{Introducció}
Text de la PAC...

\end{document}
\end{verbatim}

\subsection{Memòria de projecte (PR)}

Per a una memòria de projecte (treball més extens) acostuma a ser
convenient:
\begin{itemize}
  \item utilitzar el perfil \texttt{PR} (portada completa i sumari);
  \item personalitzar amb més cura la portada (grau, títol de la memòria,
        codi d'assignatura, etc.);
  \item establir de manera explícita les metadades PDF.
\end{itemize}

Un pont de partida podria ser:
\begin{verbatim}
\documentclass[
  PR,
  data = {
    student       = {Nom Cognom},
    grade         = {Grau en Enginyeria Informàtica},
    subj-fullname = {Nom complet de l'assignatura},
    subj-shortname= {NCA},
    subj-code     = {01312},
    ca-fullname   = {Memòria de Projecte},
    ca-shortname  = {PR}
  }
]{P3CTeX}

\begin{document}

\PECTeXconfig{
  cover-config = {
    text-backtitle = {Grau en Enginyeria Informàtica},
    text-title     = {Memòria de Projecte},
    text-minititle = {PR},
    text-author    = {Nom Cognom}
  },
  page-config = {
    author-name = {Nom Cognom},
    doc-name    = {Memòria PR},
    subj-name   = {NCA},
    refcode     = {01312}
  },
  pdf-config = {
    pdf-title  = {Memòria de Projecte - Nom de l'assignatura},
    pdf-author = {Nom Cognom}
  }
}

% Portada i sumari impresos automàticament en iniciar el document.

\sectionuoc{Introducció}
Text de la memòria...

\end{document}
\end{verbatim}

\section{Exemple complet}

L'exemple següent il·lustra un ús típic per a una PAC en català:

\begin{verbatim}
\documentclass[
  PAC,
  data = {
    student       = {Nom Cognom},
    grade         = {Enginyeria Informàtica},
    subj-fullname = {Nom complet de l'assignatura},
    subj-shortname= {NCA},
    subj-code     = {01312},
    ca-fullname   = {PAC 1: Introducció a P3CTeX},
    ca-shortname  = {PAC1}
  }
]{P3CTeX}

\begin{document}

\PECTeXconfig{
  cover-config = {
    text-minititle = {PAC1},
    text-author    = {Nom Cognom}
  },
  page-config = {
    doc-name   = {PAC 1},
    subj-name  = {NCA},
    refcode    = {01312}
  },
  pdf-config = {
    pdf-title  = {PAC 1 - Introducció a P3CTeX},
    pdf-author = {Nom Cognom}
  }
}

% Portada i sumari s'imprimeixen automàticament segons les opcions.

\sectionuoc{Introducció}
Text d'exemple...

\end{document}
\end{verbatim}

\end{document}

